%
% Chapter 7 - Future Work
%
\chapter{Future Work} \label{cap:future-work}

This chapter collects the parts of the original project scope that were not delivered within the project timeline, together with the concrete next step for moving the application from its current local/Vercel setup (Chapter~\ref{cap:deployment}) into institutional use. Each of these is scoped as a follow-up to the work already completed, rather than a redesign of it: the underlying architecture (Chapter~\ref{cap:architecture}) was built to accommodate all three without further structural change.

\section{Schedule Export (RF05)} \label{sec:future-export}

RF05 (Section~\ref{sec:functional-requirements}) requires the system to generate an academic calendar in Excel format for institutional distribution; no such export exists in the current codebase. A spreadsheet suits institutional distribution better than a static rendering of the calendar UI would: recipients can sort, filter, and re-organize a schedule of events and occurrences in a way a fixed page layout does not allow. The natural shape for this is one worksheet per academic year, with one row per occurrence and a column for each piece of metadata already modeled on the event and occurrence (name, category, classification, status, responsible, description, and dates), built with a library such as \texttt{exceljs}, which produces a native \texttt{.xlsx} workbook directly from TypeScript objects. This would follow the same Controller-Service-Data shape as the rest of the backend (Section~\ref{sec:backend-layers}): a new service method assembling the relevant events and occurrences for a given academic year and category selection, and a controller action streaming the resulting workbook back as the HTTP response.

\section{Email Notification Templates (remaining RF04 scope)} \label{sec:future-notification-templates}

RF04 originally called for configurable email or other text templates associated with events, for sending proactive notifications. The in-app half of this requirement was delivered (Section~\ref{sec:backend-notifications}): EduCal now surfaces upcoming occurrences to signed-in users through the notification bell (Section~\ref{sec:frontend-notifications}), based on a configurable per-event or per-occurrence lead time. What remains is the template half of the original requirement: a free-text field, attached to an event, that an editor could fill in with the wording of a reminder (for example, the exact phrasing to send to a department ahead of an enrolment deadline), and a way to surface that text so a staff member could copy it and send it manually through the institution's own email system; it remains future work purely because it was not reached within the project timeline, not because of any technical obstacle identified during development.

\section{Deployment to the ISEL Virtual Machine} \label{sec:future-deployment}

The application currently runs from a local development environment and a Vercel preview deployment backed by the Redis-based data layer (Section~\ref{sec:current-infra}); PostgreSQL, the storage backend the institutional deployment is meant to use, has only been exercised locally. Chapter~\ref{cap:deployment} already lays out the intended production setup, a single virtual machine on ISEL's internal network hosting both the Next.js application and PostgreSQL, and the concrete steps to provision it (Section~\ref{sec:deployment-roadmap}). None of those steps have been carried out yet, because of timeline constraints but, the immediate next steps are running the Drizzle migrations against the new database, seeding its reference lookup tables and initial administrator account, and putting the application behind a process manager and reverse proxy so it survives restarts and serves traffic over HTTPS. Establishing a recurring PostgreSQL backup schedule at that point, rather than after the calendar is already in active use, would close the last gap between the current delivery and a production-ready deployment.
